\documentclass[11pt,a4paper]{article}
\usepackage[utf8]{inputenc}
\usepackage[T1]{fontenc}
\usepackage[italian]{babel}
\usepackage{lmodern}
\usepackage[margin=2.4cm]{geometry}
\usepackage{microtype}
\usepackage{booktabs,tabularx,array}
\usepackage{enumitem}
\usepackage{xcolor}
\usepackage{xurl}
\usepackage{hyperref}
\definecolor{accent}{HTML}{315A52}
\hypersetup{colorlinks=true,linkcolor=accent,urlcolor=accent,pdftitle={Come hanno lavorato gli agenti nell'ultima run PathMNIST}}
\setlist[itemize]{topsep=4pt,itemsep=2pt}
\setlist[enumerate]{topsep=4pt,itemsep=3pt}
\setlength{\parindent}{0pt}
\setlength{\parskip}{6pt}
\newcommand{\artefatto}[1]{\nolinkurl{#1}}
\newcommand{\file}[1]{\nolinkurl{#1}}
\newcommand{\ruolo}[1]{\paragraph{#1}}

\title{Come hanno lavorato gli agenti\\nell'ultima run PathMNIST}
\author{Documento esplicativo del workflow HAT-MedMNIST}
\date{Run del 2 ottobre 2026\\Analisi degli artefatti e dei log di esecuzione}

\begin{document}
\maketitle

\begin{center}
\small\file{v2/runs/pathmnist_20261002T083530607044Z}
\end{center}

\begin{abstract}
Il documento ricostruisce il lavoro degli agenti nella run indicata, seguendo
il percorso dai dati al dossier finale. Per ciascun ruolo distingue informazioni
disponibili, decisioni, artefatti prodotti e controllo del Reviewer. La ricerca
è stata eseguita sul seed 42; la configurazione vincente è stata congelata e
allenata nuovamente sui seed 47 e 72. Il risultato è un workflow completo con
accuracy media del 92,64\%, ma con warning aperti e comportamento OOD instabile.
Completamento della pipeline e accettazione del progetto sono quindi due esiti distinti.
\end{abstract}

\tableofcontents
\newpage

\section{Il principio di collaborazione}

Gli agenti non si passano direttamente oggetti attraverso chiamate reciproche.
Il punto di scambio è un \textbf{Blackboard}: ogni agente legge gli artefatti
necessari e pubblica i propri risultati. Gli artefatti hanno un tipo, una
versione, un produttore e un checksum. I dati pesanti, come immagini,
checkpoint e predizioni, sono conservati in file referenziati dagli artefatti.

L'\textbf{Orchestrator} stabilisce l'ordine di esecuzione. Dopo ogni fase invoca
il \textbf{Reviewer}, che restituisce un rapporto di anomalie. Le decisioni e
i passaggi sono registrati in \file{decision_log.jsonl}; richieste e risposte
del reasoning sono conservate in \file{blobs/reasoning/}.

\textbf{Ibrido} significa che il calcolo di dati, training e metriche è eseguito
dal codice, mentre alcuni agenti usano il modello locale \artefatto{qwen3.8}
per proporre o motivare decisioni. Il nome del modello è registrato; il digest
dei suoi pesi è assente. Non è quindi possibile ricostruire la sua identità
esatta dal solo manifest.

\subsection{Le condizioni della run}

\begin{tabularx}{\textwidth}{@{}lX@{}}
\toprule
\textbf{Aspetto} & \textbf{Configurazione effettiva}\\
\midrule
Dataset & PathMNIST, nove classi, immagini RGB $28\times28$\\
Split & 89.996 train; 10.004 validation; 7.180 test\\
Ripetizioni & Seed 42, 47 e 72; dataset completo\\
Ricerca & Otto trial, tre round, solo sul seed 42\\
Budget dei round & 3, 6 e 9 epoche; training finale fino a 15 epoche\\
Reasoning & Ollama locale, modello richiesto; chiamate sequenziali\\
Letteratura & Estratti integrati; ricerca online disattivata\\
Ablation & Non eseguite: \artefatto{ablation_suite=false}\\
\bottomrule
\end{tabularx}

\section{La sequenza operativa}

Sul seed 42 il flusso osservato è:
\begin{enumerate}
\item Ingestion: carica il dataset e registra gli split.
\item Profiling / Ambiguity: descrive i dati e i rischi di ambiguità.
\item Prior-Art Scout: costruisce un brief di ipotesi con fonti.
\item Preprocessing: sceglie una prima rappresentazione.
\item Data Audit: verifica i dati prima del training.
\item Architecture Research: formula indicazioni per la ricerca.
\item Model Search: allena e confronta candidati usando la validation.
\item Training: riallena la configurazione selezionata.
\item Evaluation: misura il risultato sul test.
\item Abstention / OOD: calibra l'astensione e valuta le corruzioni.
\item Reporting: raccoglie il dossier, poi aggiornato dopo il confronto baseline.
\end{enumerate}

Il Reviewer interviene fra gli stadi. Nei seed 47 e 72, il punto 7 è sostituito
da \textbf{Frozen Configuration}: riprende il vincitore senza una nuova ricerca.
Le fasi iniziali di profiling e reasoning vengono comunque eseguite; non
possono cambiare la configurazione finale congelata.

\section{Gli agenti che preparano il problema}

\subsection{Ingestion: rendere disponibili dati e riferimenti}
\ruolo{Informazioni e azione}
Riceve dataset, seed e opzioni di caricamento. Carica PathMNIST preservando
gli split ufficiali e registra forma, classi, dimensioni e riferimenti degli
indici selezionati. In questa run non è attivo alcun sottocampionamento.

\ruolo{Consegna al Blackboard}
Pubblica \artefatto{dataset_manifest} e \artefatto{blob_raw_dataset}, con
il riferimento a \file{blobs/raw_dataset_v001.npz}. Il manifest permette
agli agenti successivi di conoscere il problema senza ricostruire il caricamento.

\ruolo{Esito del controllo}
Il Reviewer approva l'ingestion nei tre seed. Caricare e registrare il test
non significa renderne disponibili le metriche alla ricerca: il suo uso
valutativo è separato dalle decisioni iniziali.

\subsection{Profiling / Ambiguity: descrivere dati e incertezze}
\ruolo{Informazioni e azione}
Usa train e validation per calcolare distribuzione delle classi e statistiche
delle immagini. Il rapporto di sbilanciamento train è circa 1,63; le
distribuzioni train/validation sono molto simili. I canali RGB hanno medie
circa $(0,741;0,533;0,706)$ e deviazioni standard $(0,124;0,177;0,124)$.

L'LLM interpreta queste informazioni e segnala sovrapposizioni plausibili:
adipose/background, smooth muscle/stroma, debris e frammenti cellulari,
oltre a patch che possono contenere tessuti diversi. Queste sono
\textbf{ipotesi di ambiguità}, non errori di annotazione dimostrati.

\ruolo{Consegna al Blackboard}
Produce \artefatto{data_profile}, con statistiche, rischi e provenienza del
reasoning. La valutazione delle ambiguità non usa il risultato del test.
Il campo test imbalance presente nel profilo non va interpretato come una
misura effettiva: il test non è stato profilato in questa fase.

\subsection{Prior-Art Scout: trasformare fonti in ipotesi}
\ruolo{Informazioni e azione}
Legge il contesto del dataset e tre estratti integrati relativi a ViT,
Compact Transformers e MedMNIST v2. In questa run non esegue una ricerca web.
Le fonti hanno origine \artefatto{curated_excerpt}; il primo brief è
\artefatto{bundled}, e i successivi utilizzano la cache.

Sul seed 42 propone quattro idee: patch transformer, tokenizzazione
convoluzionale compatta, scelta della granularità delle patch e training
sensibile allo sbilanciamento. Le etichetta come ipotesi da validare.
Propone anche una composizione di augmentation e una specifica di compact
transformer. Sono proposte, non codice generato ed eseguito.

\ruolo{Consegna al Blackboard e gate}
Pubblica \artefatto{prior_art_brief}, snapshot delle fonti,
\artefatto{extension_gate_report} e \artefatto{approved_extensions}.
Le nuove proposte restano \artefatto{pending_review}: non diventano
automaticamente nuove famiglie eseguibili. Le primitive già integrate
restano utilizzabili senza approvare una nuova estensione.

Il compact transformer provato dalla ricerca appartiene al registro esistente:
\textbf{non dimostra la promozione della nuova specifica dello Scout}.
Il Reviewer segnala correttamente che estratti citabili non dimostrano
retrieval online né tutte le inferenze scientifiche sviluppate a partire da essi.

\section{Gli agenti che progettano l'esperimento}

\subsection{Preprocessing: proporre la rappresentazione iniziale}
\ruolo{Informazioni e decisione}
Legge il profilo e le ipotesi bibliografiche. Sceglie standardizzazione
train-only e una prima combinazione di augmentation:
\artefatto{hflip}, \artefatto{brightness}, \artefatto{rotate90}.
Mantiene immagini RGB e risoluzione del benchmark.

\ruolo{Consegna e aggiornamento successivo}
Pubblica la prima versione di \artefatto{representation_plan},
\artefatto{split_manifest} e del blob dei dati preparati.
Questa decisione è un punto di partenza: la ricerca può proporre una
combinazione diversa all'interno dei limiti ammessi. Il vincitore sostituisce
\artefatto{rotate90} con \artefatto{rotate180}; il Blackboard conserva
entrambe le versioni e rende esplicito il cambiamento.

\subsection{Data Audit: autorizzare l'accesso al training}
\ruolo{Informazioni e controlli}
Il controllo è deterministico: verifica split ufficiali, indici, immagini,
label, forme, valori e statistiche train-only; controlla la disgiunzione
train/validation. Nei tre seed l'audit passa e non trova sovrapposizioni.

\ruolo{Consegna e significato}
Produce \artefatto{data_audit_report}. Model Search e Training richiedono
un audit passato prima di iniziare. Un warning del Reviewer su documentazione
o interpretazione può restare aperto anche quando l'audit dei dati è positivo:
non sono lo stesso controllo.

\subsection{Architecture Research: formulare indicazioni, non un vincitore}
\ruolo{Informazioni e azione}
Combina profilo, rappresentazione e prior art in un documento di indicazioni
su famiglie, tokenizzazione, regolarizzazione e rischi. Sul seed 42 il
reasoning considera promettente il compact transformer.

\ruolo{Consegna e limite}
Pubblica \artefatto{architecture_research}. Questo artefatto guida le proposte
del search agent, ma \textbf{non seleziona il modello finale}. L'esito della
ricerca sarà una ResNet18: la preferenza teorica non prevale sulle misure
di validation. Le raccomandazioni bibliografiche sono inoltre fallibili e
non costituiscono una prova sperimentale delle relative affermazioni.

\section{La decisione centrale: ricerca e congelamento}

\subsection{Model Search sul seed 42}
L'agente legge metadati di training, indicazioni precedenti e risultati
di validation già prodotti. L'LLM propone candidati; il codice impone
schema, limiti, copertura delle famiglie e budget. La ricerca è quindi
una collaborazione fra proposte motivate e vincoli deterministici.

\begin{table}[htbp]
\centering\small
\begin{tabularx}{\textwidth}{@{}Xrr@{}}
\toprule
\textbf{Trial} & \textbf{Epoche} & \textbf{Accuracy validation}\\
\midrule
Compact transformer & 3 & 94,5522\%\\
Residual CNN bilanciata & 3 & 97,1811\%\\
ResNet18 conservativa & 3 & 97,8908\%\\
Tiny CNN regolarizzata & 6 & 92,0832\%\\
Vision transformer & 6 & 96,1515\%\\
ResNet18 conservativa & 6 & 99,0804\%\\
ResNet18 conservativa & 9 & 99,1104\%\\
ResNet18 più larga, class weighting & 9 & \textbf{99,4002\%}\\
\bottomrule
\end{tabularx}
\caption{Trial effettivamente completati. I budget diversi non consentono
un confronto a parità di costo fra famiglie.}
\end{table}

I trial sono training separati: riprovare una configurazione con più epoche
non implica proseguire automaticamente dal precedente checkpoint.
La scelta considera i trial al budget massimo; usa accuracy entro la
tolleranza dichiarata e macro-F1 per il ranking. Le metriche del test
sono escluse. Il vincitore è \artefatto{resnet18_wider_cw}.

L'agente pubblica \artefatto{search_plan}, artefatti dei singoli trial,
\artefatto{search_report}, \artefatto{train_config} e
\artefatto{best_configuration}; aggiorna anche la rappresentazione definitiva
e salva \file{best_config_v001.yaml}.

\subsection{Frozen Configuration sui seed 47 e 72}
L'agente riprende il \artefatto{best_configuration} del seed 42 e cambia
il seed di training. Ricostruisce i dati preparati e il file YAML locale,
senza fare altri trial. Questo è il meccanismo che rende confrontabili
le tre ripetizioni del modello selezionato.

Il valore di validation 99,4002\% riportato nella selezione dei tre seed
rimane quello del trial originale. Le nuove validation finali appartengono
invece ai rispettivi training. I file YAML locali hanno hash diversi
perché cambia il seed; gli altri parametri sono uguali.

\subsection{Quale agente non è intervenuto}
\artefatto{ExperimentDesignAgent} è disponibile per il percorso senza
ricerca iterativa, ma \textbf{non è stato eseguito in questa run}. La scelta
è stata affidata a Model Search, poi ripresa da Frozen Configuration.
Le ablation della rappresentazione, pur implementate come procedura,
non sono state attivate.

\section{Gli agenti che eseguono e misurano}

\subsection{Training: eseguire la scelta congelata}
Legge dati preparati, audit e configurazione finale. Allena il modello con
Lightning fino al budget di 15 epoche; registra la storia delle metriche
e conserva il checkpoint scelto sulla validation. Non propone una nuova
architettura né adatta la scelta al test.

Produce \artefatto{train_result} e il riferimento del modello.
La migliore epoca è 15, 15 e 14 nei tre seed; tutti completano 15 epoche.
Le migliori accuracy di validation finali sono rispettivamente
99,4702\%, 99,6002\% e 99,5602\%.

\subsection{Evaluation: aprire il test per la misura finale}
Usa il checkpoint selezionato e produce probabilità, logits, predizioni
e metriche del test: accuracy, macro-F1, balanced accuracy, AUC,
matrice di confusione e risultati per classe. Pubblica
\artefatto{prediction_artifact} e \artefatto{evaluation_report}.

\begin{center}
\begin{tabular}{@{}lrrr@{}}
\toprule
\textbf{Seed} & \textbf{Accuracy} & \textbf{Macro-F1} & \textbf{Balanced accuracy}\\
\midrule
42 & 92,2981\% & 0,898925 & 0,901332\\
47 & 92,1448\% & 0,898564 & 0,897181\\
72 & 93,4680\% & 0,917116 & 0,915932\\
\bottomrule
\end{tabular}
\end{center}

Il test resta fuori dalla ricerca, ma il suo risultato viene successivamente
letto per valutazione, revisione e reporting. Il divario validation/test
rimane di circa 6--7,5 punti. Lo stroma associato al tumore ha recall
fra 56,53\% e 64,37\%: il risultato aggregato non elimina questa debolezza.

\subsection{Abstention / OOD: decidere quando rinviare}
L'agente confronta max-softmax ed entropia predittiva usando la validation
e corruzioni gaussiane controllate. Sceglie lo score con migliore AUROC
medio sulla validation e calibra una soglia per coprirne circa l'80\%.
Applica poi score e soglia al test senza cambiarli in base al risultato.

Produce \artefatto{abstention_report}, curve rischio/copertura, probabilità
delle corruzioni e \artefatto{human_review_queue}. La coda registra indici
dei campioni, predizione, confidenza e motivo del rinvio.

\begin{center}\small
\begin{tabular}{@{}llrrr@{}}
\toprule
\textbf{Seed} & \textbf{Score} & \textbf{Coverage} & \textbf{Acc. accettati} & \textbf{Rinvii}\\
\midrule
42 & Entropia & 67,37\% & 97,42\% & 2.343\\
47 & Max-softmax & 65,29\% & 97,57\% & 2.492\\
72 & Max-softmax & 68,70\% & 97,67\% & 2.247\\
\bottomrule
\end{tabular}
\end{center}

L'astensione funziona sui campioni puliti nel senso di ridurre l'errore
fra quelli accettati. Non garantisce però protezione OOD: sul seed 47
l'AUROC medio è 0,3938 e alla corruzione forte è circa 0,0018. Il modello
predice background con alta confidenza e accetta quasi tutti i casi corrotti.
Il problema è già osservabile sulla validation; il selettore sceglie
il migliore fra due score insufficienti e non restituisce un esito
esplicito ``nessuno score idoneo''.

Il Reviewer segnala questo fallimento. Il seed 72 passa il criterio AUROC
minimo implementato, ma accetta circa metà dei casi con rumore forte.
Queste misure riguardano un proxy di corruzione, non OOD clinico generale.
La coda dimostra un handoff predisposto: non una revisione umana conclusa.

\section{Reviewer e Orchestrator: controllo e gestione dei warning}

Il Reviewer combina invarianti deterministiche con un esame LLM degli
artefatti disponibili. Produce rapporti \artefatto{anomaly_*} e distingue
\artefatto{continue}, \artefatto{revise} e \artefatto{stop}.
Un finding critico deterministico può bloccare il flusso; le affermazioni
LLM non sufficientemente dimostrate rimangono advisory.

\artefatto{revise} non equivale a una modifica già eseguita. L'Orchestrator
può ripetere lo stadio se è disponibile una remediation sicura; altrimenti
registra \artefatto{review_unresolved} e promuove lo stadio con warning.
Nei log compaiono 5, 6 e 2 eventi di questo tipo sui seed 42, 47 e 72.
La run non mostra un ciclo di correzione che risolva tutti i rilievi.

Fra i rilievi osservati figurano fonti integrate invece di retrieval web,
inferenze che eccedono le citazioni, debolezza dello stroma, digest LLM
assente e fallimento OOD del seed 47. Alcuni commenti sulle proposte di
augmentation richiedono interpretazione: una ricetta nuova pending non
impedisce di usare singolarmente primitive già integrate.

\textbf{Il Reviewer segnala e rende tracciabile il problema; non è un agente
che ripara autonomamente il modello.} L'Orchestrator permette il completamento
operativo con warning. La procedura di accettazione finale mantiene però
questi problemi aperti.

\section{Reporting, baseline e accettazione}

Reporting assembla \file{dossier.json}, registro degli artefatti e stato
del workflow. Dopo la prima chiusura, il runner esegue una baseline
convenzionale sullo stesso dataset e seed, pubblica
\artefatto{baseline_report} e \artefatto{comparison_report}, richiama
il Reviewer e aggiorna il dossier.

\textbf{La baseline è una procedura del runner, non un ulteriore agente
di reasoning.} L'accuracy media è 82,99\%, contro 92,64\% della pipeline.
Il vantaggio di 9,65 punti riguarda però un sistema con 6,29 milioni di
parametri contro una Tiny CNN di 5.433 parametri, oltre a rappresentazione
e costi diversi. Non misura isolatamente il contributo degli agenti.

La procedura di accettazione associa requisiti a evidenze e produce
\artefatto{acceptance_report}. Tutti i seed risultano tecnicamente completi,
ma non pronti per l'accettazione e senza evidenza TRL~7 completa.
Passa ora il requisito di ripetizione con configurazione congelata.
Restano assenti ablation, prova di retrieval online, fault evidence sul
codice esatto, replay completo, mitigazioni approvate e sign-off indipendente.

\section{Che cosa chiarisce questa run sul lavoro degli agenti}

La collaborazione osservata è coerente con la separazione fra proposta,
esecuzione e controllo. Gli agenti iniziali descrivono il problema;
quelli di reasoning suggeriscono rappresentazioni e candidati; il search
produce evidenze di validation; training ed evaluation eseguono e misurano;
il Reviewer rende visibili limiti e incongruenze; reporting mantiene
il dossier finale.

La decisione architetturale più importante è stata presa sul seed 42 e
conservata nei due successivi. L'accuracy migliora rispetto alla run
precedente, ma l'instabilità OOD mostra che congelare gli iperparametri
non basta a garantire stabilità di confidenza e rifiuto.

Questa run documenta quindi \textbf{un workflow agentico operativo e auditabile},
con decisioni effettive dell'LLM e risultati misurati. Non dimostra ancora
che tutte le motivazioni siano scientificamente corrette, che l'OOD sia
affidabile o che l'orchestrazione sia superiore a una ricerca convenzionale
a parità di spazio e budget.

\appendix
\section{Mappa delle evidenze}

I riferimenti seguenti sono relativi alla directory della run indicata
in apertura. Gli artefatti versionati sono in \file{seed_*/artefacts/}.

\begin{tabularx}{\textwidth}{@{}lX@{}}
\toprule
\textbf{Domanda} & \textbf{Evidenza principale}\\
\midrule
Quali opzioni sono state usate? & \artefatto{run_configuration}, \artefatto{run_manifest}\\
Che cosa è avvenuto e in quale ordine? & \file{seed_*/decision_log.jsonl}\\
Che cosa ha detto l'LLM? & \file{seed_*/blobs/reasoning/}\\
Quali dati e statistiche? & \artefatto{dataset_manifest}, \artefatto{data_profile}\\
Quali fonti e proposte? & \artefatto{prior_art_brief}, \artefatto{extension_gate_report}\\
Quali scelte di input? & Versioni di \artefatto{representation_plan}\\
Quali candidati e vincitore? & Artefatti \artefatto{trial_*}, \artefatto{search_report}, \artefatto{best_configuration}\\
Quale training finale? & \artefatto{train_result}, \file{best_config_v001.yaml}\\
Quali risultati e rinvii? & \artefatto{evaluation_report}, \artefatto{abstention_report}, \artefatto{human_review_queue}\\
Quali rilievi restano? & Rapporti \artefatto{anomaly_*}, \file{acceptance_report.json}\\
Quale riepilogo fra seed? & \file{experiment_summary.json}\\
\bottomrule
\end{tabularx}

La ricostruzione si basa sugli artefatti e sui log della run. Il codice corrente
è stato consultato per interpretare i meccanismi, ma ha un hash diverso dal
codice registrato nella run: non sostituisce lo snapshot originale.
Le verifiche successive sono conservate in
\file{v2/audits/pathmnist_20261002T083530607044Z/}; non sono eventi eseguiti dagli
agenti durante la run. Il presente documento non modifica gli artefatti storici.

\end{document}
